\documentclass[11pt,a4paper,twoside]{report}

% Paquetes esenciales
\usepackage[utf8]{inputenc}
\usepackage[spanish,es-tabla]{babel}
\usepackage[left=2.5cm,right=2.5cm,top=3cm,bottom=3cm]{geometry}
\usepackage{graphicx}
\usepackage{float}
\usepackage{booktabs}
\usepackage{tabularx}
\usepackage{amsmath}
\usepackage{amsfonts}
\usepackage{amssymb}
\usepackage{xcolor}
\usepackage{fancyhdr}
\usepackage{titlesec}
\usepackage{microtype}
\usepackage[hidelinks]{hyperref}
\usepackage{listings}
\usepackage{natbib}
\usepackage[T1]{fontenc}
\usepackage{lmodern} % Usa la fuente Latin Modern, que es vectorial
% Configuración de colores
\definecolor{nycred}{RGB}{255,90,95} % Color brand Airbnb
\definecolor{codegray}{rgb}{0.95,0.95,0.95}
\definecolor{codeblue}{rgb}{0.1,0.2,0.8}

% Configuración de listados de código
\lstset{
    backgroundcolor=\color{codegray},
    commentstyle=\color{gray},
    keywordstyle=\color{codeblue}\bfseries,
    numberstyle=\tiny\color{gray},
    stringstyle=\color{nycred},
    basicstyle=\ttfamily\footnotesize,
    breakatwhitespace=false,         
    breaklines=true,                 
    captionpos=b,                    
    keepspaces=true,                 
    numbers=left,                    
    numbersep=5pt,                  
    showspaces=false,                
    showstringspaces=false,
    showtabs=false,                  
    tabsize=2
}

% Configuración de cabeceras y pies de página
\pagestyle{fancy}
\fancyhf{}
\fancyhead[LE,RO]{\thepage}
\fancyhead[RE]{\leftmark}
\fancyhead[LO]{\rightmark}
\renewcommand{\headrulewidth}{0.4pt}

% Título y Autor
\title{
    \vspace{-2cm}
    \Huge \textbf{Predicción y Segmentación Dinámica del Mercado Inmobiliario en NYC}\\[0.5cm]
    \Large Informe Técnico Completo del Proyecto de Machine Learning
}
\author{
    \textbf{Autor:} Sergio Minguez Cruces\\
    \textbf{Institución:} UNIE Universidad\\
    \textbf{Grado:} Matemáticas\\
    \textbf{Asignatura:} Minería de Datos
}
\date{14 de Mayo de 2026}

\begin{document}

% Portada
\maketitle

% Resumen Ejecutivo
\begin{abstract}
Este informe detalla el diseño, implementación y evaluación de un pipeline de \textit{Machine Learning} end-to-end sobre el ecosistema de alquileres de Airbnb en la ciudad de Nueva York (2019). El proyecto consolida un análisis exhaustivo de casi 50,000 propiedades, abordando problemas de asimetría de datos, valores faltantes condicionales (MAR) y multicolinealidad. Mediante ingeniería de variables avanzadas (distancias geoespaciales), técnicas de reducción de dimensionalidad (PCA), y un ecosistema de modelización supervisada robusto (ensembles y SVM), se logró un modelo altamente generalizable ($R^2 \approx 59\%$) para la predicción de precios. De manera complementaria, el aprendizaje no supervisado (K-Means) ha permitido segmentar la oferta en nichos de mercado claramente definidos. Todo el sistema está orquestado con una adherencia estricta a principios de reproducibilidad y culmina en el despliegue de una herramienta interactiva.
\end{abstract}

% Índices
\tableofcontents
\listoffigures
\listoftables

\clearpage

% Inclusión de capítulos
\chapter{Introducción y Motivación}

\section{Contextualización del Proyecto}
La ciudad de Nueva York representa uno de los ecosistemas inmobiliarios y turísticos más densos, dinámicos y competitivos a nivel mundial. Con millones de visitantes anuales, plataformas de alojamiento compartido como Airbnb han revolucionado la manera en que se distribuye y consume el espacio habitable temporal. El conjunto de datos analizado en este trabajo, \textit{New York City Airbnb Open Data} de 2019 \cite{airbnb2019}, captura una radiografía transversal de este fenómeno, conteniendo información detallada de casi 50,000 listados activos.

\section{Identificación del Problema de Negocio}
Para una plataforma bilateral como Airbnb, mantener el equilibrio entre la oferta (anfitriones) y la demanda (huéspedes) es fundamental para la maximización de ingresos. Un problema persistente en este ecosistema es la **correcta valoración del precio nocturno** de las propiedades. Los anfitriones novatos a menudo carecen del conocimiento del mercado para fijar precios competitivos, lo que resulta en tasas de ocupación subóptimas (si el precio es excesivo) o pérdida de ganancias potenciales (si el precio es deficiente).

Este proyecto aborda la necesidad de crear un motor algorítmico capaz de estimar el valor intrínseco de un alojamiento basándose estrictamente en sus características observables (tipo de habitación, métricas de disponibilidad) y, críticamente, en su centralidad geográfica y ubicación.

\section{Objetivos del Proyecto}
El desarrollo técnico de este informe se estructuró para dar cumplimiento a tres objetivos analíticos principales:

\begin{enumerate}
    \item \textbf{Predicción de Precios en NYC (Supervisado):} Construir, optimizar y contrastar una batería extensa de modelos de regresión (desde modelos lineales penalizados hasta métodos de \textit{ensemblaje} avanzado como XGBoost \cite{chen2016xgboost} y Stacking) para predecir el precio del alojamiento de forma precisa y generalizable.
    \item \textbf{Inteligencia Espacial (Ingeniería de Variables):} Trascender la simple pertenencia a un vecindario mediante el cálculo de distancias reales a puntos neurálgicos de la ciudad, transformando metadatos de latitud y longitud en variables predictivas robustas.
    \item \textbf{Segmentación de Mercado (No Supervisado):} Implementar un proceso de agrupación algorítmica para descubrir submercados latentes dentro de NYC, perfilando la oferta más allá de divisiones administrativas, lo que permite detectar patrones de demanda y concentración de opciones \textit{premium} versus económicas.
\end{enumerate}

Para garantizar un rigor estadístico absoluto y reproducibilidad, el proyecto completo (desde el remuestreo inicial hasta el \textit{GridSearchCV}) ha sido anclado mediante un estricto control estocástico (\texttt{random\_state=42}).

\chapter{Análisis Exploratorio de Datos (EDA)}

\section{Descripción General del Conjunto de Datos}
El conjunto de datos original cuenta con 48,895 registros y 16 variables que engloban tanto identificadores nominales, métricas temporales, como coordenadas geográficas continuas. Antes de someter estos datos a cualquier transformación predictiva, se realizó un Análisis Exploratorio de Datos (EDA) para diagnosticar problemas estructurales de asimetría, colinealidad y valores perdidos.

\section{Tratamiento de Valores Faltantes y Naturaleza MAR}
Una revisión de completitud reveló valores faltantes focalizados en cuatro columnas: \texttt{name}, \texttt{host\_name}, \texttt{last\_review} y \texttt{reviews\_per\_month}. 
Las dos primeras representan una pérdida negligente (menos del 0.05\%) y fueron imputadas con la etiqueta genérica "Unknown".

El caso de \texttt{last\_review} y \texttt{reviews\_per\_month} requirió un análisis diagnóstico más detallado. Se identificó que la falta de datos en estas variables no obedece a un error de recolección (MCAR, \textit{Missing Completely At Random}), sino que corresponde a una estructura MAR (\textit{Missing At Random}). Específicamente, existe una dependencia lógica determinista: estas variables son invariablemente nulas cuando la variable \texttt{number\_of\_reviews} es igual a 0. Esto significa que las propiedades que nunca han sido alquiladas o evaluadas carecen por definición de tasa de reseña mensual o fecha de última reseña. En base a esta evidencia, \texttt{reviews\_per\_month} fue imputado con 0, y \texttt{last\_review} fue transformada en una nueva variable derivada que cuantifica los días desde la última interacción, asignando un umbral máximo a los casos sin reseñas previas.

\section{Análisis de Distribuciones y Asimetría Severa}
El diagnóstico distribucional univariante alertó de desviaciones significativas respecto a la normalidad en la amplia mayoría de las variables predictoras continuas y en la variable objetivo (\texttt{price}). 

\begin{itemize}
    \item \textbf{Variable Objetivo (\texttt{price}):} La distribución del precio presenta un severo sesgo a la derecha (\textit{right skewness}), con una masa principal concentrada por debajo de los \$200 USD y valores atípicos extremos superando los \$10,000 USD. Para estabilizar la varianza y mejorar el comportamiento de los modelos lineales y paramétricos, se aplicó una transformación logarítmica $\log(1+x)$ (conocida como \texttt{log1p}).
    \item \textbf{Variables Temporales y de Actividad:} Las variables \texttt{minimum\_nights}, \texttt{number\_of\_reviews} y la derivada de recencia exhibieron un patrón de asimetría similar dictado por distribuciones de tipo Ley de Potencias. Por ende, la transformación \texttt{log1p} fue también aplicada a estos predictores.
\end{itemize}

Esta estandarización logarítmica mitiga sustancialmente la influencia desproporcionada de los valores atípicos sin necesidad de realizar recortes drásticos que suprimirían información valiosa (por ejemplo, listados genuinamente de lujo o penthouses).

\section{Distribución Espacial Inicial}
Visualmente, existe una hegemonía clara en los precios. El distrito de \textbf{Manhattan} ostenta los valores promedio más altos (aproximadamente \$196/noche), seguido de Brooklyn (\$124/noche), mientras que el Bronx y Staten Island fungen como áreas puramente residenciales y periféricas con valores significativamente menores. Asimismo, la variable \texttt{room\_type} polariza el mercado: las propiedades de tipo \textit{Entire home/apt} dominan casi con exclusividad la franja superior de precios.

\chapter{Ingeniería de Variables e Inteligencia Espacial}

\section{Creación de Features Espaciales (Haversine)}
La ubicación cruda definida por latitud y longitud, si bien provee un marco de referencia absoluto, carece de semántica de negocio si no se interrelaciona con los puntos neurálgicos de la demanda turística de Nueva York. 

Para dotar de inteligencia espacial al modelo, se instrumentó el cálculo matemático de distancias en superficie esférica utilizando la fórmula del semiverseno (Haversine). Se proyectó la distancia en kilómetros desde cada propiedad hacia tres focos principales de actividad económica y turística:
\begin{enumerate}
    \item \textbf{Times Square} (Centro de la actividad turística y comercial).
    \item \textbf{Central Park} (Atracción paisajística primaria).
    \item \textbf{Grand Central Terminal} (Eje de interconexión logística del transporte).
\end{enumerate}

Esta ingeniería transformó las coordenadas geográficas estáticas en métricas dinámicas de centralidad, convirtiéndose subsecuentemente en uno de los ejes de mayor \textit{feature importance} global.

\section{Codificación Híbrida de Variables Categóricas}
Los modelos matemáticos (especialmente los basados en cálculo de distancias como SVM y Regresión Lineal) no asimilan características textuales. La naturaleza de los datos categóricos dictó una estrategia de codificación híbrida:
\begin{itemize}
    \item \textbf{One-Hot Encoding (OHE):} Aplicado a variables con baja cardinalidad como \texttt{neighbourhood\_group} (5 valores) y \texttt{room\_type} (3 valores), evitando la introducción de una jerarquía artificial inexistente.
    \item \textbf{Target Encoding:} Aplicado al vector \texttt{neighbourhood} (más de 220 valores nominales), sustituyendo el nombre del vecindario por el precio logarítmico medio del mismo. Esto colapsó la altísima dimensionalidad sin mermar la varianza y reteniendo su peso explicativo.
\end{itemize}

\section{Reducción de Dimensionalidad: Análisis de Componentes Principales (PCA)}
La expansión del dataset mediante el \textit{Target Encoding} y las distancias espaciales, junto a las variables temporales preexistentes, incrementó sustancialmente el riesgo de multicolinealidad estructural (por ejemplo, la alta dependencia entre \texttt{availability\_365} y métricas temporales de reservas).

Para destilar la señal del ruido, se aplicó el Análisis de Componentes Principales (PCA). PCA encuentra una transformación lineal ortogonal del espacio de características originales hacia un nuevo sistema de coordenadas donde la varianza se maximiza iterativamente.

La ecuación de proyección a un nuevo componente principal $z_k$ a partir del vector estandarizado $x$ es:
\begin{equation}
z_k = \mathbf{w}_k^T \mathbf{x}
\end{equation}
Sujeto a la restricción ortogonal $||\mathbf{w}_k|| = 1$.

En nuestra implementación, conservamos un umbral de hiperparámetro fijado empíricamente en el 95\% de varianza explicada. Este abordaje permitió eliminar componentes residuales redundantes, acelerando el tiempo de convergencia (\textit{training time}) de modelos complejos como *Support Vector Machines* en la fase predictiva y estableciendo un espacio isotrópico ideal para el *Clustering*.

\chapter{Aprendizaje No Supervisado: Segmentación de Mercado}

\section{Metodología de Clustering}
Antes de acometer la tarea predictiva, resultaba imperativo comprender si la oferta de alquileres en Nueva York obedecía a distribuciones homogéneas o si, por el contrario, existían submercados latentes. Se empleó el algoritmo \textbf{K-Means} operando estrictamente sobre el espacio latente reducido (generado por PCA), garantizando así que las distancias Euclidianas calculadas por el algoritmo estuviesen exentas del sesgo impuesto por magnitudes dispares o colinealidad intrínseca.

La función de coste a minimizar (Inercia o WCSS - \textit{Within-Cluster Sum of Squares}) se define como:
\begin{equation}
J = \sum_{j=1}^{k} \sum_{i=1}^{n} ||x_i^{(j)} - c_j||^2
\end{equation}

\section{Determinación por Consenso del Número Óptimo de Clústeres (K)}
La dependencia analítica a un único método heurístico (e.g., Regla del Codo) para la selección de $K$ puede inducir sesgos visuales e interpretativos severos. Este proyecto apostó por una metodología de validación tripartita y unificada:
\begin{enumerate}
    \item \textbf{Método del Codo (Elbow Method):} Análisis del punto de máxima curvatura sobre la curva de inercia decreciente, identificando el punto de saturación donde la ganancia marginal de añadir centroides adicionales colapsa.
    \item \textbf{Análisis de Silueta (Silhouette Score):} Maximización de la cohesión intra-cluster frente a la separación inter-cluster.
    \item \textbf{Estadístico Gap (Gap Statistic):} Comparación de la dispersión de los datos frente a una distribución nula aleatoria de referencia.
\end{enumerate}
A partir de estos tres cálculos, el algoritmo extrajo dinámicamente la **mediana estadística** (\textit{Consensus K}) de las tres predicciones heurísticas, fijando la partición de la base de datos de manera agnóstica a intervenciones manuales.

\section{Perfilado y Caracterización}
La segmentación resultante arrojó perfiles de mercado sustancialmente disímiles y dotados de semántica interpretable. La inyección de inteligencia analítica sobre los resultados permitió aislar las modas estadísticas de la variable \texttt{room\_type} por cada clúster. 

Los hallazgos demuestran que el mercado neoyorquino está fuertemente fragmentado por una conjunción bivariante entre el tipo de habitación ofertada y el eje centralidad-precio. Se identificaron clústeres representativos de **ofertas premium** (habitaciones de tipo \textit{Entire home/apt} de alto costo ubicados a mínimas distancias de los hitos espaciales) frente a **clústeres residuales** (mayoritariamente dominados en porcentaje por opciones \textit{Shared room} o \textit{Private room} en las periferias del Bronx o zonas externas de Queens y Brooklyn). 

Este hallazgo es crucial: dictamina que una estrategia estandarizada de precios (\textit{One-size-fits-all}) fracasaría inevitablemente, justificando la necesidad subyacente de emplear algoritmos de regresión flexibles capaces de emular ramificaciones condicionales.

\chapter{Aprendizaje Supervisado: Modelado Predictivo}

\section{Diseño Experimental y Ecosistema de Modelos}
El objetivo principal (\textit{core}) de la investigación residía en la construcción de un motor de inferencia capaz de aproximar la valoración económica del mercado inmobiliario. 
Para asegurar el rigor, la evaluación se rigió bajo un paradigma estricto de \textit{Hold-out Test} (partición 80/20) asistido por \textbf{Cross-Validation} estructurada (K-Folds). Adicionalmente, todo el protocolo incluyó bloqueos de semillas (\texttt{RANDOM\_STATE=42}) para imposibilitar la fuga de datos o discrepancias en la reproducibilidad.

Dada la heterogeneidad y los picos de no-linealidad expuestos en la fase de clustering, se entrenó masivamente un ecosistema integrado por 12 algoritmos distintos:
\begin{enumerate}
    \item \textbf{Modelos de Regresión Lineal y Penalizados:} Regresión Lineal OLS (Baseline), Regresión Ridge (penalización $L_2$) y Regresión Lasso (penalización $L_1$).
    \item \textbf{Basados en Árboles:} Árbol de Decisión (CART).
    \item \textbf{Support Vector Machines:} LinearSVR y SVR con kernel Radial (SVR\_RBF).
    \item \textbf{Ensembles (Bagging y Boosting):} RandomForestRegressor, AdaBoost, XGBoost \cite{chen2016xgboost} y Bagging sobre regresores SVR.
    \item \textbf{Meta-Modelos (Ensamblajes de Alto Nivel):} Voting Regressor y Stacking Regressor (ensamblajes heterogéneos anidados).
\end{enumerate}

\section{Optimización Paramétrica (GridSearchCV)}
La extracción de máxima capacidad teórica a los algoritmos se orquestó a través de un proceso de optimización exhaustiva, \texttt{GridSearchCV}. El grid de búsqueda iteró sobre los principales parámetros de control de la arquitectura de red algorítmica: profundidades de árbol (\texttt{max\_depth}), márgenes duros y blandos para el caso SVM (\texttt{C} y \texttt{gamma}), y coeficientes de penalización ($\alpha$). Se usó el coeficiente de determinación ($R^2$) como métrica unificadora de convergencia en este espacio cartesiano paramétrico.

\section{Evaluación, Diagnóstico y Selección Final}
Más allá de las métricas puramente aisladas como el Error Absoluto Medio (MAE) o el $R^2$, el protocolo de auditoría incorporó una métrica penalizante diseñada en este estudio: el \textbf{\textit{Overfitting Gap}}. Esta dimensión calculaba la sustracción analítica entre la varianza explicada en la partición de entrenamiento (Train) y la de evaluación (Hold-out). Cualquier modelo superando un salto del 0.05 era clasificado como inestable (memorización de ruido o varianza de sesgo extrema) y descartado categóricamente.

\begin{table}[H]
    \centering
    \caption{Resultados del Hold-out Test para modelos estables (\textit{Gap} $< 0.05$)}
    \begin{tabular}{lcccc}
        \toprule
        \textbf{Modelo} & \textbf{Hold-out $R^2$} & \textbf{MAE (\$)} & \textbf{RMSE (\$)} & \textbf{Overfitting Gap} \\
        \midrule
        \textbf{SVR\_RBF} & \textbf{0.5898} & \textbf{54.29} & \textbf{194.85} & \textbf{-0.0026} \\
        Voting\_Assembly & 0.5866 & 55.15 & 193.95 & 0.0414 \\
        XGBoost & 0.5863 & 55.40 & 193.10 & 0.0287 \\
        LinearRegression & 0.5471 & 57.32 & 196.29 & -0.0199 \\
        \bottomrule
    \end{tabular}
\end{table}

\textbf{Modelo Ganador:} El algoritmo \textbf{Support Vector Regressor con Kernel Radial (SVR\_RBF)} se posicionó inequívocamente como la solución definitiva. Su kernel no lineal ($RBF$) probó ser la herramienta matemática perfecta para capturar las interacciones geolocalizadas complejas dentro del espacio PCA. Alcanzó un $R^2$ aproximado del 59\% con un MAE competitivo de \$54.29 dólares y un \textit{Overfitting Gap} negativo, certificando matemáticamente su sublime e intachable capacidad de generalización algorítmica ante nueva casuística.

\chapter{Despliegue e Interfaz Interactiva}

\section{Arquitectura de Producción (Streamlit)}
La finalidad de cualquier ecosistema analítico aplicado es su transferencia al dominio humano mediante una interfaz interpretable. El modelo \textbf{SVR\_RBF} ganador, la proyección topológica PCA, la matriz de escalado StandardScaler y el propio modelado de segmentación (K-Means) fueron encapsulados matemáticamente y exportados a disco (\textit{Pickle Serialization}).

Este esquema actúa como motor \textit{backend} para una interfaz de usuario completamente dinámica construida bajo la arquitectura Streamlit. La aplicación elimina la necesidad de que los \textit{stakeholders} interactúen vía terminal, ofreciendo en una arquitectura web una experiencia integral unificada, desde tableros de diagnóstico y representaciones biplot del PCA, hasta un simulador \textit{real-time}.

\section{El Simulador de Precios}
El pilar de despliegue principal es la interfaz de inferencia interactiva. Se dotó de una arquitectura isomórfica que garantiza la consistencia del vector analítico: la información digitada por el usuario (Barrio, Tipo de habitación, Disponibilidad, etc.) cruza una réplica exacta del pipeline subyacente. 
Esto incluye el cálculo algebraico \textit{on-the-fly} de las distancias Haversine (empleando matrices NumPy para minimizar latencia), su escalado y posterior proyección dimensional en PCA previo a ser transferido a la unidad SVR para su evaluación.

\section{Cartografía Geográfica Integrada (Folium)}
Para fortalecer el carácter holístico y la usabilidad (UX) del simulador, el ingreso numérico por coordenadas fue sustituido por una selección visual a través de mapas interactivos utilizando \textit{Folium}. 
El sistema superpone máscaras logísticas sobre el área del barrio (Neighbourhood) validando los \textit{inputs} del usuario antes de integrarlos al ciclo predictivo. Así, no sólo se retorna una recomendación de la rentabilidad (Predicción del Precio por Noche) sino que, paralelamente, el \textit{endpoint} indica automáticamente al grupo segmentado (\textit{Cluster}) al cual la propiedad analizada pertenecería bajo la perspectiva no supervisada.

\chapter{Conclusiones y Trabajo Futuro}

\section{Síntesis de Cumplimiento de Objetivos}
Las metas delineadas en la fundamentación de este informe, específicamente alineadas con los requisitos del examen final estandarizado, fueron cumplimentadas en su totalidad de manera holística:
\begin{enumerate}
    \item \textbf{Resolución Predictiva:} Se logró modelar la alta varianza del mercado de Nueva York (donde dominaban patrones de ruido y asimetría temporal severa), absorbiendo un 59\% de la varianza total a través de un algoritmo robusto a sobreajustes (\textit{SVR\_RBF}).
    \item \textbf{Ingeniería Híbrida:} La superación del Baseline se fundamentó mayoritariamente por la codificación logarítmica para lidiar con el MAR detectado y la adición del preprocesamiento de coordenadas Haversine, reducidas vectorialmente en un espacio estable vía PCA.
    \item \textbf{Conocimiento Latente:} El *Clustering* dinámico brindó la validación final que las interacciones del mercado de alquiler están determinadas sustancialmente por la exclusividad de uso y la distancia a los focos masivos de consumo (\textit{Times Square}, etc.).
\end{enumerate}

\section{Limitaciones del Estudio}
El rigor científico exige el reconocimiento formal de los límites arquitectónicos presentes:
\begin{itemize}
    \item \textbf{Falta de Variabilidad Estacional:} Al operar exclusivamente sobre una instantánea fotográfica del año 2019, es matemáticamente imposible extraer factores de picos de demanda estival, eventos temporales masivos o la subsecuente pandemia global del COVID-19, dictando que el modelo captura estructuras estructurales puras, pero no variaciones macroeconómicas transitorias.
    \item \textbf{Carencia de Contexto Visual o Descriptivo:} En mercados altamente subjetivos, factores no estructurados como la limpieza general inferida por las fotografías de la galería y la verborragia del *Host* en su descripción textual son predictores ocultos inalcanzables en los datos tabulares disponibles.
\end{itemize}

\section{Recomendaciones y Trabajo Futuro}
Con base en los resultados, el \textit{Action Plan} para futuras iteraciones y para un posible salto hacia la comercialización de la herramienta recae en:
\begin{enumerate}
    \item \textbf{Análisis de Sentimiento Computacional (NLP):} Procesar masivamente el corpus textual de los comentarios empleando modelos de Lenguaje Neuronal masivos (LLMs), agregando un coeficiente de penalización/bonificación basado en polaridad.
    \item \textbf{Motor de Retroalimentación en Tiempo Real:} Diseñar la arquitectura analítica orientada a flujos y *Event Streams*, re-entrenando pesos dinámicos en los modelos SVR con APIs para adaptar las rentabilidades predichas a shocks de la demanda de alta inmediatez logística.
\end{enumerate}


% Bibliografía
\bibliographystyle{apalike}
\bibliography{referencias}

\end{document}
